%!tex

\magnification=1400

\input tikz
\usetikzlibrary{arrows.meta, math}
\input ../../../TeX/macros/macros.tex

\def\ct#1{\hfil#1\hfil}
\def\tm{\,(\hbox{s})}

\center{%
  {DEPARTAMENTO DE F\'ISICA}
}
\vskip -.5ex
\center{%
  {UNIVERSIDADE FEDERAL DA PARA\'IBA}
}
\vskip -.5ex
\center{%
  {F\'ISICA EXPERIMENTAL I}
}
\vskip 2ex
\center{%
  {\sl Experimento 5\quad \rm P\^endulo simples}
}
\vskip 2ex
\par

\noindent
\vbox{
  \halign{
    #\hfil &\quad # \cr
    {\bsl Docente} & {\bsl Turma} \cr
     & F\'ISICA EXPERIMENTAL I\cr
  }
}

\noindent{\bsl Grupo}\par
\indent
\vbox{
  \halign{
   # \hfil &\quad # \cr
  {\bsl Nome} & {\bsl Matr\'\i cula}\cr
  Harllen Ara\'ujo de Sena & \cr
  }
}

\procl{Quest\~ao}{}{%
  Explique porque escolhemos medir um intervalo de tempo $t$, equivalente a v\'arios per\'\i odos $T$, em vez de medir diretamente um \'unico per\'\i odo de oscila\c c\~ao $T$.
}

Medir uma \'unica oscila\c c\~ao exigiria muita precis\~ao. Por este motivo, mede-se um intervalo de tempo correspondente a um n\'umero conveniente de oscila\c c\~oes. Por exemplo, digamos que em 10.81 segundos contamos 10 idas e voltas, isto \'e oscila\c c\~oes. O per\'\i odo \'e precisamente $10.81\,\hbox{s}/10=0.1081\,\hbox{s}$.


\procl{Quest\~ao}{}{%
  Para cada comprimento $L$, realize 5 medidas do tempo total $t$ consecutivas do p\^endulo e anote os resultados na {\bsl Tabela 1} abaixo.
}

\bigskip
\center{%
  \vbox{
    \halign{
     \ct{#} & \ct{#} & \ct{#} & \ct{#} & \ct{#} & \ct{#} & \ct{#} \cr 
     $L\,(\hbox{m})$ & $t_1\tm$ & $t_2\tm$ & $t_3\tm$ & $t_4\tm$ & $t_5\tm$ & $\bar{t}\tm$\cr
           &&&&& \cr    
     0.350 & 11.62 & 11.78 & 11.50 & 11.85 & 11.76 & 11.70 \cr    
     0.496 & 13.70 & 14.35 & 14.20 & 14.17 & 10.03 & 13.29 \cr
     0.648 & 16.15 & 16.19 & 15.72 & 15.87 & 16.19 & 16.02 \cr
     0.783 & 17.96 & 17.67 & 17.64 & 17.47 & 17.88 & 17.72 \cr
     0.899 & 18.72 & 19.20 & 19.10 & 18.83 & 19.05 & 18.98 \cr
     1.097 & 20.63 & 20.62 & 21.17 & 20.82 & 20.43 & 20.73 \cr
     1.193 & 21.61 & 22.04 & 21.47 & 21.96 & 21.48 & 21.71 \cr
     1.373 & 23.05 & 23.07 & 23.66 & 23.47 & 23.87 & 23.42 \cr   
    }
  }
}
\vskip 1em
\center{%
  {\bf Tabela 1}
}
\bigskip

\procl{Quest\~ao}{}{%
  A partir dos resultados obtidos na {\bsl Tabela 1}\/, determine o valor m\'edio do per\'\i odo $\bar{T}$, para cada comprimento $L$ do fio. Calcule $\bar{T}^2$ e insira todos os resultados na {\bsl Tabela 2} abaixo.
}

\bigskip
\center{%
  \vbox{
    \halign{
     \ct{#} & \ct{#} & \ct{#} \cr
     $L\,(\hbox{m})$ & $\bar{T}\tm$ & $\bar{T}^2\tm$ \cr
           &       &       \cr    
     0.350 & 1.170 & 1.369 \cr    
     0.496 & 1.329 & 1.766 \cr
     0.648 & 1.602 & 2.568 \cr
     0.783 & 1.772 & 3.141 \cr
     0.899 & 1.898 & 3.602 \cr
     1.097 & 2.073 & 4.299 \cr
     1.193 & 2.171 & 4.714 \cr
     1.373 & 2.342 & 5.487 \cr   
    }
  }
}
\vskip 1em
\center{%
  {\bf Tabela 1}
}
\bigskip
 

\procl{Quest\~ao}{}{%
  Fa\c ca o gr\'afico de $L$ versus $\bar{T}$. Determine incialmente qual grandeza ficar\'a no eixo horizontal. H\'a uma depend\^encia linear?
}

A grandeza sendo mensurada \'e o per\'\i odo, que por sua vez depende do comprimento do fio do p\^endulo.

\fig{%
  \tikzpicture[scale=1.2]
    \def\cn#1#2{
      \node [
        pin={
          [fill=white, draw=black]
          #2:{$#1$}
        },
        circle,
        draw=black, 
        fill=white,
        inner sep=1.5pt,
      ] at #1 {};
    }
    \def\err(#1,#2){
      \def\errl{.05}
      \def\errt{.01}
      \draw [draw=red] (#1,#2) circle [radius=1pt];
      \draw (#1-\errl,#2) -- (#1+\errl,#2);
      \draw (#1,#2-\errt) -- (#1,#2+\errt);
    }
    \coordinate (p1) at (0.350,1.170);
    \coordinate (p2) at (0.496,1.329);
    \coordinate (p3) at (0.648,1.602);
    \coordinate (p4) at (0.783,1.772);
    \coordinate (p5) at (0.899,1.898);
    \coordinate (p6) at (1.097,2.073);
    \coordinate (p7) at (1.193,2.171);
    \coordinate (p8) at (1.373,2.342);
    \scope
      \cn{(0.350,1.170)}{235}
      \cn{(0.496,1.329)}{300}
      \cn{(0.648,1.602)}{345}
      \cn{(0.783,1.772)}{0}
      \cn{(0.899,1.898)}{180}
      \cn{(1.097,2.073)}{160}
      \cn{(1.193,2.171)}{3}
      \cn{(1.373,2.342)}{110}
    \endscope
    \scope[xshift=4cm, scale=1.5]
      \draw[->] (-.25,0) -- (2,0) node [above=2pt] {L};
      \draw[->] (0,-.25) -- (0,3) node [above=2pt] {$\bar{T}$};
      \err(0.350,1.170)
      \err(0.496,1.329)
      \err(0.648,1.602)
      \err(0.783,1.772)
      \err(0.899,1.898)
      \err(1.097,2.073)
      \err(1.193,2.171)
      \err(1.373,2.342)
    \endscope
  \endtikzpicture
}{%
  gr\'afico $\bar{T}$ versus $L$.
}

N\~ao. A depend\^encia \'e dada por teoria via aproxima\c c\~oes por \^angulos pequenos por
$$
  T=2\pi\sqrt{L\over g},
$$
uma rela\c c\~ao quadr\'atica.

\procl{Quest\~ao}{}{%
  Fa\c ca um segundo gr\'afico de $L$ e $\bar{T}$. Agora h\'a uma depend\^encia linear? O modelo te\'orico prev\^e esta depend\^encia? Prove.
}

Antes, observemos que para \^angulos pequenos vale 
$$
  T=2\pi\sqrt{L\over g},
$$
da\'\i\ vem
$$
  T^2={4\pi^2\over g}L,
$$
e portanto a dep\^endencia \'e linear. Em conformidade, a teoria prev\^e esta dep\^endencia.

Segue o gr\'afico requerido.

\fig{%
  \tikzpicture[scale=2]
    \def\err(#1,#2){
      \def\errl{.05}
      \def\errt{.01}
      \draw (#1-\errl,#2) -- (#1+\errl,#2);
      \draw (#1,#2-\errt) -- (#1,#2+\errt);
    }
    \draw[->] (-.25,0) -- (2,0) node [above=2pt] {L};
    \draw[->] (0,-.25) -- (0,6) node [above=2pt] {$\bar{T}^2$};
    \err(0.350,1.369)
    \err(0.496,1.766)
    \err(0.648,2.568)
    \err(0.783,3.141)
    \err(0.899,3.602)
    \err(1.097,4.299)
    \err(1.193,4.714)
    \err(1.373,5.487)
  \endtikzpicture
}{%
  gr\'afico $\bar{T}^2$ versus $L$.
}

\procl{Quest\~ao}{}{%
  A partir do m\'etodo dos m\'\i nimos quadrados obtenha o coeficiente angular e sua incerteza. Apresente todos os c\'alculos.
}

 Antes, fa\c camos uma an\'alise grosseira do procedimento de determina\c c\~ao dos coeficientes $A$ e $B$, tais que
$$
  (\Lambda)\qquad y_i=A+Bx_i
$$
em que $x_i$ e $y_i$ s\~ao as vari\'aveis aleat\'orias (medidas) em an\'alise. Segundo J. R. Taylor a medida de $y_i$ \'e governada pela distribui\c c\~ao Gaussiana, com o mesmo comprimento $\sigma_y$ para todas as medidas.

Em sua obra nota-se que a rela\c c\~ao $(\Lambda)$, implicar\'a que a probabilidade de se obter o valor $y_i$ \'e
$$
  {\rm P}_{AB}(y_i)\propto{1\over\sigma_y}\exp\biggl({(A+Bx_i-y_i)^2\over2\sigma_y^2}\biggr)
$$
e a probabilidade de se obter o conjunto completo de medidas $y_1,\ldots,y_n$ \'e o produto

$$
  (\Pi)\qquad P_{AB}(y_1,\ldots,y_n)\;=\;\prod_{i=1}^nP_{AB}(y_i)\;\propto\;{1\over\sigma_y^n}\exp\Bigl(-{\chi^2\over2}\Bigr)
$$
em que evidentemente
$$
  \chi^2={1\over\sigma_y^2}\sum_{i=1}^n(A+Bx_i-y_i)^2
$$
Assume-se ent\~ao, que as melhores estimativas para as inc\'ognitas $A$ e $B$ concernentes as medidas $x_i$ e $y_i$, s\~ao tais que $P_{AB}(y_1,\ldots,y_n)$ \'e m\'axima. Em virtude de $(\Pi)$ podemos inferir que isto \'e alca\c cado quando $\chi^2$ \'e m\'\i nimo. Sabemos que $\chi^2$ \'e uma fun\c c\~ao de $2$ vari\'aveis reais, nomeadamente $A$ e $B$ com valores reais. Ademais,

$$
  \nabla\chi^2(A,B)=-{2\over\sigma_y^2}\Bigl(\sum_iA+Bx_i-y_i,x_i\sum_iA+Bx_i-y_i\Bigr)
$$
Seu \'unico ponto cr\'\i tico \'e determinado fazendo-se $\nabla\chi^2=0$, equivalentemente solucionar o sistema
$$
  (S)\qquad
  \left\{
  \eqalign{
    nA+S_xB   &=S_y\cr
    S_xA+Q_xB &=P\cr
  }
  \right.
$$
em que
$$
  S_x=\sum_ix_i,\quad S_y=\sum_iy_i,\quad Q_x=\sum_ix_i^2\quad \hbox{e}\quad P=\sum_ix_iy_i.
$$
Resolvendo para $A$ e $B$, temos
$$
  A={S_yQ_x-PS_x\over nQ_x-S_x^2}\quad\hbox{e}\quad B={nP-S_xS_y\over nQ_x-S_x^2}.
$$
Por fim, notemos que
$$
 [(\nabla\chi^2)']=
 \left[\matrix{
  \partial_A^2\chi^2  & \partial_{AB}\chi^2 \cr
  \partial_{BA}\chi^2 & \partial_B^2\chi^2 \cr
 }
 \right]
 =
 \left[\matrix{
  \displaystyle{2n\over\sigma_y^2}   & \displaystyle{2S_x\over\sigma_y^2} \cr
  \omit\cr
  \displaystyle{2S_x\over\sigma_y^2} & \displaystyle{2Q_x\over\sigma_y^2} \cr
 }
 \right],
$$
cuja forma quadr\'atica Hessiana \'e
$$
  \det[(\nabla\chi^2)']={4\over\sigma_y^4}(nQ_x-S_x^2).
$$
Sejam $x=(x_1,\ldots,x_n)$ e $u=(1,\ldots,1)$, pela desigualdade de {\bsl Cauchy-Bunyakovski-Schwartz} inferimos
$$
  S_x=\langle x,u\rangle\leq\sqrt{\sum_ix_i^2}\sqrt{n}=\sqrt{Q_x}\sqrt{n},
$$
donde conclu\'\i mos que
$$
  nQ_x-S_x^2\geq0.
$$
Al\'em disso, o m\'etodo dos m\'\i nimos quadrados s\'o \'e significativo se n\~ao for o caso que todos os $x_i$ s\~ao iguais. Sob est\'a hip\'otesea, segue-se que
$$
  \det[(\nabla\chi^2)']={4\over\sigma_y^4}(nQ_x-S_x^2)>0.
$$
Portanto, o ponto cr\'\i tico $(A,B)$ \'e ponto de m\'\i nimo local de $\chi^2$.

Em nosso caso
$$
  \vbox{
    \halign{
      \hfil$#$\hfil\cr
      x=(0.350, 0.496, 0.648, 0.783, 0.899, 1.097, 1.193, 1.373);\cr
      y=(1.369, 1.766, 2.568, 3.141, 3.602, 4.299, 4.714, 5.487).\cr
    }
  }
$$
Agora, relembremos que
$$
  A={S_yQ_x-PS_x\over nQ_x-S_x^2}\quad\hbox{e}\quad B={nP-S_xS_y\over nQ_x-S_x^2}.
$$
Pondo estes dados num pequeno script em python, obtemos:
$$
  A=-0.10431716936864718\quad\hbox{e}\quad B=4.062227324893873
$$
Segundo J. R. Taylor, uma boa estimativa para $\sigma_y$ \'e dada de
$$
  (\Pi)\qquad P_{AB}(y_1,\ldots,y_n)\;\propto\;{1\over\sigma_y^n}\exp\Bigl(-{\chi^2\over2}\Bigr)
$$
pelo princ\'\i pio de probabilidade m\'axima. Fazendo
$$
  S=\sum_{i=1}^n(A+Bx_i-y_i)^2
$$
e derivando $(\Pi)$ com respeito a $\sigma_y$, e solucionando para zero, i.e.,
$$
  -{n\over\sigma_y^{n+1}}\exp\Bigl(-{\chi^2\over2}\Bigr)+{1\over\sigma_y^n}\exp\Bigl(-{\chi^2\over2}\Bigr){S\over\sigma_y^3}=0
$$
temos
$$
  \sigma_y=\sqrt{S\over n}=\sqrt{{1\over n}\sum_{i=1}^n(A+Bx_i-y_i)^2}.
$$
Por quest\~oes de compensa\c c\~ao no c\'alculo de $A$ e $B$ aumenta-se $\sigma_y$ ({\it vide} {\bsl An introduction to Error Analysis}\/), fazendo-se a seguinte adapta\c c\~ao em $\sigma_y$:
$$
  \sigma_y=\sqrt{{1\over n-2}\sum_{i=1}^n(A+Bx_i-y_i)^2}.
$$
Por quest\~oes de escassez de tempo, referindo-se novamente J. R. Taylor tem-se que
$$
  \sigma_B=\sigma_y\sqrt{n\over nQ_x-S_x^2}.
$$
Novamente, usando um script python temos
$$
  \sigma_B=0.0829044567347054
$$
Finalmente, aproximando os resultados de $B$ e $\sigma_B$ para o primeiro d\'\i gito significativo, obtemos
$$
  B=4.06\pm0.08.
$$

\procl{Quest\~ao}{}{%
  Determine o valor da acelera\c c\~ao da gravidade $g$ na superf\'{\i}cie da Terra a partir do coefiente do termo dependente encontrado no item anterior, e, sua incerteza usando a teoria de propaga\c c\~ao de erros. Forne\c ca o resultado na forma $g\pm\Delta g$. 
}

Primeiramente, temos da teoria que
$$
  T^2=\Bigl({4\pi^2\over g}\Bigr)L.
$$
Comparando-se com nossa determina\c c\~ao estat\'\i stica obtemos
$$
  B={4\pi^2\over g},
$$
i.e.,
$$
  g={4\pi^2\over B},
$$
sabendo que $B\approx4.06$, temos
$$
  g=9.723748178413162.
$$
Pela teoria de propraga\c c\~ao de erros tem-se
$$
  \delta_g=\sqrt{({\partial_Bg}\cdot\sigma_y)^2}=\sigma_B|\partial_Bg|=\sigma_B\Bigl({2\pi\over B}\Bigr)^2.
$$
Agora, substituindo as aproxima\c c\~oes de $\sigma_B$ e $B$ determinadas na quest\~ao anterior incorremos naturalmente
$$
  \delta_g=0.19160094932833815
$$
aproximando para o primeiro d\'\i gito siginificativo obtemos que $\delta_g=0.19$. Por fim,
$$
  g\pm\delta_g\approx9.72\pm0.19.
$$

\procl{Quest\~ao}{}{%
  Compare a precis\~ao do seu resultado com o valor esperado $g=9.78\pm0.01\;{\rm m}/{\rm s}^2$. Os valores s\~ao compat\'\i veis entre si?
}

Note que nosso experimento deu um intervalo $[9.53,9.91]$ que intersecta o intervalo esperado $[9.77,9.79]$, na realidade o experimental inclui o estipulado.

\fig{%
  \tikzpicture[scale=2]
    \def\err(#1,#2){
      \def\errl{.05}
      \def\errt{.01}
      \def\mult{1}
      \draw [Bracket-Bracket](#1-\mult*\errl,#2) -- (#1+\mult*\errl,#2);
      \draw (#1,#2-\mult*\errt) -- (#1,#2+\mult*\errt);
    }
    \err(0.350,1.369)
    \err(0.496,1.766)
    \err(0.648,2.568)
    \err(0.783,3.141)
    \err(0.899,3.602)
    \err(1.097,4.299)
    \err(1.193,4.714)
    \err(1.373,5.487)
    \draw  (0.223,0.8015595240826864) -- (1.5,5.989023817972162);
  \endtikzpicture
}{%
  Regress\~ao de $\bar T^2$ em $L$. 
}

\bye

